% ================================================================
% The printable version of latex/starters/targeted/KS4/geometry/vectors/vectorJourneysStarter_retrieveAndConnect.tex
%
% This is the same deck, laid out to be handed out rather than put on
% the board. Every slide is shown as it stands before any answer is
% revealed, and 4 copies of it go on one page of A4, two by two, so
% that a printed page cuts into four question sheets.
%
% Every slide of the deck is here, the title page included, and each takes
% exactly one page. So page 1 is the first slide, page 2 the second, and
% printing the pages you want is how you print the starters you want.
%
%   made from           latex/starters/targeted/KS4/geometry/vectors/vectorJourneysStarter_retrieveAndConnect.tex
%   pages on the sheet  2, one per slide of the deck
%   copies of each      4, two by two on A4 landscape
%   printing rules      version 3
%
% Nothing was reworded, nothing was rewritten and nothing was left out.
% Each slide is the slide from the deck, asked for at its first step,
% which is how the answers stay hidden.
%
% Please don't edit this file. It is written again from the one it was
% made from whenever that changes, so an edit here would be lost - edit
% that file instead and this one follows.
% ================================================================
% ================================================================
% The Retrieve and Connect version of latex/starters/targeted/KS4/geometry/vectors/vectorJourneysStarter.tex
%
% This is the same deck as the one above, made by renaming the titles in it.
% Every slide that was a numbered starter is now titled
% "Retrieve and Connect", and the word is renamed wherever else a title
% used it. Nothing outside a title has been changed.
%
%   made from           latex/starters/targeted/KS4/geometry/vectors/vectorJourneysStarter.tex
%   retitling rules     version 1
%
% Please don't edit this file. It is written again from the one it was
% made from whenever that changes, so an edit here would be lost - edit
% that file instead and this one follows.
% ================================================================
\documentclass{beamer}
\usepackage{tikz}
\usetikzlibrary{arrows.meta}
\title{Year 10 Vector Retrieve and Connects}
\date{}

\usepackage{multicol}

% ================================================================
% Answer-overlay helpers
% ================================================================
\newcommand{\ablank}[1]{%
  \alt<2>{\textcolor{red}{#1}}{\underline{\phantom{#1}}}%
}
\newcommand{\ashow}[1]{\uncover<2->{\textcolor{red}{\small #1}}}
\newcommand{\ashowq}[1]{\alt<2->{\textcolor{red}{\small #1}}{?}}

% Answers drawn straight on to a TikZ diagram, revealed on overlay 2 like \ashow.
% Space is reserved on overlay 1, so the diagram does not move between slides.
\usetikzlibrary{overlay-beamer-styles}
\tikzset{
  ans/.style     = {red, thick, visible on=<2->},
  ansfill/.style = {red!15, visible on=<2->},
  anslab/.style  = {red, font=\tiny, visible on=<2->},
}

% ---------------------------------------------------------------
% Added so this deck prints four slides to a page. Nothing above
% this line was changed.
% ---------------------------------------------------------------
\usepackage{pgfpages}

% 4 slides to a sheet of A4, two by two, with a thin line round each
% so that the page can be cut up.
%
% Each slide keeps a margin of 1mm around it and no more, so the four fill
% the sheet and the pieces they cut into are as big as the paper allows. Only
% one direction actually comes out that tight: a slide is a different shape
% from a quarter of a sheet of A4, so it is scaled to fit and the slack it
% cannot use is left along whichever pair of edges it did not fill.
\pgfpagesuselayout{4 on 1}[a4paper,landscape,border shrink=1mm]
\pgfpageslogicalpageoptions{1}{border code=\pgfusepath{stroke}}
\pgfpageslogicalpageoptions{2}{border code=\pgfusepath{stroke}}
\pgfpageslogicalpageoptions{3}{border code=\pgfusepath{stroke}}
\pgfpageslogicalpageoptions{4}{border code=\pgfusepath{stroke}}

% There is nothing on paper to click, and a slide announcing a new section
% would sit between two copies of the same question and put every page after
% it out of step with the grid.
\setbeamertemplate{navigation symbols}{}
\AtBeginPart{}
\AtBeginSection{}
\AtBeginSubsection{}
\AtBeginSubsubsection{}
\begin{document}

\frame<1>{\titlepage}
\frame<1>{\titlepage}
\frame<1>{\titlepage}
\frame<1>{\titlepage}

%------------------------------------------------
\begin{frame}<1>{Retrieve and Connect}

\begin{multicols}{2}
\begin{enumerate}
  \item Start with zero, add 7, subtract 10, add 6, subtract 3, what is the result? \ashow{0}
  \item What if you started with $-5$? \ashow{$-5$}
  \item Add:
  $\begin{pmatrix}4\\1\end{pmatrix}
  +
  \begin{pmatrix}2\\3\end{pmatrix}$ \ashow{\(\binom{6}{4}\)}
  \item Find:
  $-\begin{pmatrix}5\\2\end{pmatrix}$ \ashow{\(\binom{-5}{-2}\)}
  \item What column vector represents the reverse of the following vector:
  $\begin{pmatrix}3\\-4\end{pmatrix}$? \ashow{\(\binom{-3}{4}\)}
\end{enumerate}
\end{multicols}

\vspace{0.3cm}



\begin{columns}[T]

\column{0.45\textwidth}
\begin{center}
\begin{tikzpicture}[scale=0.7]

\coordinate (A) at (0,0);
\coordinate (C) at (2,1);
\coordinate (B) at (4,2);

\draw[->] (A)--(C) node[midway,above] {$\mathbf{a}$};
\draw[->] (C)--(B) node[midway,above] {$\mathbf{b}$};

\draw[->,dashed] (A)--(B);

\node[below left] at (A) {A};
\node[above right] at (B) {B};

\end{tikzpicture}
\end{center}

\column{0.55\textwidth}

\begin{enumerate}
    \setcounter{enumi}{5}
  \item Write an expression in terms of $\bf{a}$ and $\bf{b}$ that represents the whole journey from A to B. \ashow{\(\mathbf{a}+\mathbf{b}\)}
  \item Which vector would undo that journey? How many different ways can you write it? \ashow{\(\overrightarrow{BA} = -(\mathbf{a}+\mathbf{b}) = -\mathbf{a}-\mathbf{b}\)}
\end{enumerate}

\end{columns}

\end{frame}
\begin{frame}<1>{Retrieve and Connect}

\begin{multicols}{2}
\begin{enumerate}
  \item Start with zero, add 7, subtract 10, add 6, subtract 3, what is the result? \ashow{0}
  \item What if you started with $-5$? \ashow{$-5$}
  \item Add:
  $\begin{pmatrix}4\\1\end{pmatrix}
  +
  \begin{pmatrix}2\\3\end{pmatrix}$ \ashow{\(\binom{6}{4}\)}
  \item Find:
  $-\begin{pmatrix}5\\2\end{pmatrix}$ \ashow{\(\binom{-5}{-2}\)}
  \item What column vector represents the reverse of the following vector:
  $\begin{pmatrix}3\\-4\end{pmatrix}$? \ashow{\(\binom{-3}{4}\)}
\end{enumerate}
\end{multicols}

\vspace{0.3cm}



\begin{columns}[T]

\column{0.45\textwidth}
\begin{center}
\begin{tikzpicture}[scale=0.7]

\coordinate (A) at (0,0);
\coordinate (C) at (2,1);
\coordinate (B) at (4,2);

\draw[->] (A)--(C) node[midway,above] {$\mathbf{a}$};
\draw[->] (C)--(B) node[midway,above] {$\mathbf{b}$};

\draw[->,dashed] (A)--(B);

\node[below left] at (A) {A};
\node[above right] at (B) {B};

\end{tikzpicture}
\end{center}

\column{0.55\textwidth}

\begin{enumerate}
    \setcounter{enumi}{5}
  \item Write an expression in terms of $\bf{a}$ and $\bf{b}$ that represents the whole journey from A to B. \ashow{\(\mathbf{a}+\mathbf{b}\)}
  \item Which vector would undo that journey? How many different ways can you write it? \ashow{\(\overrightarrow{BA} = -(\mathbf{a}+\mathbf{b}) = -\mathbf{a}-\mathbf{b}\)}
\end{enumerate}

\end{columns}

\end{frame}
\begin{frame}<1>{Retrieve and Connect}

\begin{multicols}{2}
\begin{enumerate}
  \item Start with zero, add 7, subtract 10, add 6, subtract 3, what is the result? \ashow{0}
  \item What if you started with $-5$? \ashow{$-5$}
  \item Add:
  $\begin{pmatrix}4\\1\end{pmatrix}
  +
  \begin{pmatrix}2\\3\end{pmatrix}$ \ashow{\(\binom{6}{4}\)}
  \item Find:
  $-\begin{pmatrix}5\\2\end{pmatrix}$ \ashow{\(\binom{-5}{-2}\)}
  \item What column vector represents the reverse of the following vector:
  $\begin{pmatrix}3\\-4\end{pmatrix}$? \ashow{\(\binom{-3}{4}\)}
\end{enumerate}
\end{multicols}

\vspace{0.3cm}



\begin{columns}[T]

\column{0.45\textwidth}
\begin{center}
\begin{tikzpicture}[scale=0.7]

\coordinate (A) at (0,0);
\coordinate (C) at (2,1);
\coordinate (B) at (4,2);

\draw[->] (A)--(C) node[midway,above] {$\mathbf{a}$};
\draw[->] (C)--(B) node[midway,above] {$\mathbf{b}$};

\draw[->,dashed] (A)--(B);

\node[below left] at (A) {A};
\node[above right] at (B) {B};

\end{tikzpicture}
\end{center}

\column{0.55\textwidth}

\begin{enumerate}
    \setcounter{enumi}{5}
  \item Write an expression in terms of $\bf{a}$ and $\bf{b}$ that represents the whole journey from A to B. \ashow{\(\mathbf{a}+\mathbf{b}\)}
  \item Which vector would undo that journey? How many different ways can you write it? \ashow{\(\overrightarrow{BA} = -(\mathbf{a}+\mathbf{b}) = -\mathbf{a}-\mathbf{b}\)}
\end{enumerate}

\end{columns}

\end{frame}
\begin{frame}<1>{Retrieve and Connect}

\begin{multicols}{2}
\begin{enumerate}
  \item Start with zero, add 7, subtract 10, add 6, subtract 3, what is the result? \ashow{0}
  \item What if you started with $-5$? \ashow{$-5$}
  \item Add:
  $\begin{pmatrix}4\\1\end{pmatrix}
  +
  \begin{pmatrix}2\\3\end{pmatrix}$ \ashow{\(\binom{6}{4}\)}
  \item Find:
  $-\begin{pmatrix}5\\2\end{pmatrix}$ \ashow{\(\binom{-5}{-2}\)}
  \item What column vector represents the reverse of the following vector:
  $\begin{pmatrix}3\\-4\end{pmatrix}$? \ashow{\(\binom{-3}{4}\)}
\end{enumerate}
\end{multicols}

\vspace{0.3cm}



\begin{columns}[T]

\column{0.45\textwidth}
\begin{center}
\begin{tikzpicture}[scale=0.7]

\coordinate (A) at (0,0);
\coordinate (C) at (2,1);
\coordinate (B) at (4,2);

\draw[->] (A)--(C) node[midway,above] {$\mathbf{a}$};
\draw[->] (C)--(B) node[midway,above] {$\mathbf{b}$};

\draw[->,dashed] (A)--(B);

\node[below left] at (A) {A};
\node[above right] at (B) {B};

\end{tikzpicture}
\end{center}

\column{0.55\textwidth}

\begin{enumerate}
    \setcounter{enumi}{5}
  \item Write an expression in terms of $\bf{a}$ and $\bf{b}$ that represents the whole journey from A to B. \ashow{\(\mathbf{a}+\mathbf{b}\)}
  \item Which vector would undo that journey? How many different ways can you write it? \ashow{\(\overrightarrow{BA} = -(\mathbf{a}+\mathbf{b}) = -\mathbf{a}-\mathbf{b}\)}
\end{enumerate}

\end{columns}

\end{frame}

%------------------------------------------------

\end{document}